% !TeX encoding = UTF-8
\documentclass[]{article}

\usepackage{imakeidx}
\makeindex[]

\usepackage[]{graphicx, svg, float}

\title{\vspace{-4cm}\includesvg[scale=.15]{figures/rustacean-flat-gesture}\\ \sffamily\bfseries{}Rust Programming}
\author{\sffamily\bfseries{}Skylar Coelho}
\date{\sffamily\bfseries\today}

\usepackage{sectsty}

% Sets all sectioning headers (Title, Section, Subsection, etc.) to Sans-Serif
\allsectionsfont{\sffamily\bfseries}

\usepackage[]{minted}

\usepackage[]{pagecolor}

\pagecolor{darkgray}
\definecolor{softwhite}{RGB}{235, 219, 178}
\color{softwhite}

\usepackage[]{hyperref}
\hypersetup{colorlinks=true, linkcolor=yellow}

\newcommand{\code}[1]{\textbf{\texttt{\color{orange}#1}}}

\begin{document}

\maketitle

\section{Rust Basics}
\subsection{Installing Rust}
Install Rust from the \href{https://www.rust-lang.org/tools/install}{Rust website}.
Once you have rust installed, check that the Rust compiler is set up correctly by running
\begin{minted}[linenos, style=gruvbox-dark, bgcolor=darkgray]{console}
	$ rustc --version
\end{minted}
Rust can be updated using rustup, the included tool chain for installing and maintaining Rust.
The version of rustup can be checked in a similar matter.
Rust can be updated via rustup by running 
\begin{minted}[linenos, style=gruvbox-dark, bgcolor=darkgray]{console}
	$ rustup --version
	$ rustup update
\end{minted}
\subsection{Cargo Basics}
Rust runs typically from a main function created by initializing the project using cargo.
Initialize a project using cargo
\begin{minted}[linenos, style=gruvbox-dark, bgcolor=darkgray]{console}
	$ cargo new project_name
\end{minted}
If the directory already exists, you can build a project in that directory using
\begin{minted}[linenos, style=gruvbox-dark, bgcolor=darkgray]{console}
	$ cargo init
\end{minted}
You can add dependencies, like external libraries from \href{crates.io}{crates.io}, by running
\begin{minted}[linenos, style=gruvbox-dark, bgcolor=darkgray]{console}
	$ cargo add package_name
\end{minted}
Finally, when you're ready, you can build your project for debugging or release using cargo as well via
\begin{minted}[linenos, style=gruvbox-dark, bgcolor=darkgray]{console}
	$ cargo build # compiles for debug by default
	$ cargo build --release # for building using the release profile
	$ cargo run # compiles for debug and runs the executable
	$ cargo run --release # compiles and runs the release profile executable
\end{minted}
Additionally, you can always compile your Rust script using the Rust compiler itself and then running the produced executable
\begin{minted}[linenos, style=gruvbox-dark, bgcolor=darkgray]{console}
	$ rustc main.rs # compile
	$ ./program_name # run the produced executable
\end{minted}
\subsection{Basic Setup}
Typical Rust programs will contain the following code snippet within \texttt{main.rs} by default
\begin{minted}[linenos, style=gruvbox-dark, bgcolor=darkgray]{rust}
	fn main(){
		println!("Hello, world!")
	}
\end{minted}
The code contained in the main function is what gets compiled and called by the typical Rust compilation actions.
Note that the default function setup contains the macro \texttt{println!} which prints the string literal \texttt{"Hello, world!"} to the console.
We know this is a macro because the standard naming scheme for macros in Rust is to leave an "!" at the end of a function name.
\subsection{Including Libraries}
Certain libraries are packaged with Rust, known as the \textit{standard library}, shortened to \textit{std}.
These packages do not need to be included as dependencies via cargo and can be brought into your code via the following syntax
\begin{minted}[linenos, style=gruvbox-dark, bgcolor=darkgray]{rust}
	use std::io
\end{minted}
The double colon '::' is similar to the dot '.' notation from other languages where the above brings in the \texttt{io} library from the standard library.
\end{document}